\section{Why a fixed Iwasawa degree does not bound the weight sum}
\label{sec:fixed-iwasawa-growth-obstruction}

AC-K26-size allows the specialized index to be arbitrarily large
when the power series varies.  The following stronger algebraic
test keeps one polynomial fixed and tests every positive integer
weight.  It isolates the quantitative input still required in
AC-K27; it is not a counterexample involving arithmetic elliptic
units or Galois cohomology.

\begin{theorem}[AC-K28: a fixed degree-two specialization obstruction]
\label{AC-K28}
Let $p$ be an odd prime, and let
$B:\mathbb N_{\ge1}\to\mathbb R_{\ge0}$ be any finite-valued
nondecreasing function.  There is a monic polynomial
$f\in\mathbb Z_p[T]$ of degree two with the following properties.
It is distinguished, has $\mu=0$ and Weierstrass degree two,
and
\[
 f((1+p)^e-1)\ne0\qquad(e=1,2,\ldots).
                                                        \tag{IW.1}
\]
Nevertheless, for every $N\ge1$,
\[
 \sum_{\substack{1\le e\le N\\e\equiv1\bmod p-1}}
       v_p\bigl(f((1+p)^e-1)\bigr)\ \ge B(N).
                                                        \tag{IW.2}
\]
Thus the failure is possible even on one fixed Teichm\"uller
weight progression.  The coefficients of $f$ are allowed to be
arbitrary $p$-adic integers; no assertion of their algebraicity
over $\mathbb Q$ or of a bounded arithmetic height is made.
\end{theorem}
\begin{proof}
\proofstep{1. Two interlaced sequences of integer approximants.}
Set $c=p-1$, $e_1=1$, and $e_2=p$.  Suppose $e_j,e_{j+1}$
and the earlier integers $k_i$ have been chosen.
Choose an integer $k_j\ge2$ such that
\[
 k_j\ge B(e_{j+1}),\qquad
 k_j\ge k_{j-1}+2\quad(j>1),\qquad
 e_j+c p^{k_j}>e_{j+1},
\]
and define $e_{j+2}=e_j+c p^{k_j}$.
These are only lower bounds on $k_j$, so the recursion continues
indefinitely.  The sequence $e_j$ is strictly increasing and
unbounded.  Every $e_j$ is congruent to one modulo $p-1$.
Define in $\mathbb Z_p$
\[
 \alpha=1+c\sum_{j\ {
m odd}}p^{k_j},\qquad
 \beta=p+c\sum_{j\ {
m even}}p^{k_j}.             \tag{IW.3}
\]
The increasing exponents make both series convergent.
For odd $j$, $e_j$ is the partial sum defining $\alpha$
before exponent $k_j$; for even $j$ it is the corresponding
partial sum for $\beta$.  The first omitted summand has
valuation exactly $k_j$, and every later summand has larger
valuation.  Consequently
\[
 v_p(e_j-\alpha)=k_j\quad(j\ {
m odd}),\qquad
 v_p(e_j-\beta)=k_j\quad(j\ {
m even}).          \tag{IW.4}
\]
Neither $\alpha$ nor $\beta$ is an ordinary integer.
Indeed, their $p$-adic digit expansions have infinitely many
digits $p-1$ at the displayed, separated exponents, and
infinitely many zero digits between them.  The initial digits
at exponents zero or one do not overlap these terms, since
$k_j\ge2$.  A nonnegative integer has eventually zero digits;
a negative integer has eventually $p-1$ digits.  Both
possibilities are excluded.  The two roots are also distinct,
since $\alpha\equiv1$ and $\beta\equiv0\pmod p$.

\proofstep{2. The cyclotomic change of variable.}
For $x,y\in\mathbb Z_p$ one has
\[
 v_p\bigl((1+p)^x-(1+p)^y\bigr)
                         =1+v_p(x-y)\quad(x\ne y).
                                                        \tag{IW.5}
\]
Here $(1+p)^x$ is defined by continuous extension from the
ordinary integer powers.  For completeness, if $p\nmid m$,
the binomial expansion gives
$v_p((1+p)^m-1)=1$.  If $z\in p\mathbb Z_p$ is nonzero,
the term $pz$ in $(1+z)^p-1$ has strictly smaller valuation
than every other term, since $p$ is odd; hence raising
$1+z$ to the $p$th power raises that valuation by one.
Induction gives the formula $1+v_p(m)$ for every nonzero
integer $m$; negative $m$ follows by taking a unit inverse.
It both proves existence of the continuous extension and,
by approximating $x,y$ by integers with their difference's
valuation stabilized, proves (IW.5).

Put $a=(1+p)^\alpha-1$, $b=(1+p)^\beta-1$, and
$f(T)=(T-a)(T-b)$.  Both $a,b$ belong to $p\mathbb Z_p$.
Thus $f$ is monic of degree two and its two lower
coefficients belong to $p\mathbb Z_p$, as asserted.
Equation (IW.5) makes the cyclotomic map injective;
since neither exponent is an integer, (IW.1) follows.
For every positive integer $e$, the same equation gives
the finite nonnegative valuation
\[
 v_p\bigl(f((1+p)^e-1)\bigr)
            =2+v_p(e-\alpha)+v_p(e-\beta).          \tag{IW.6}
\]
For $e=e_j$, one of the last two terms is $k_j$ by
(IW.4); the other is zero because the residues of the
two roots modulo $p$ differ.

\proofstep{3. Every cutoff is covered.}
For any $N\ge1$ there is a unique $j$ with
$e_j\le N<e_{j+1}$.  The term $e=e_j$ occurs in (IW.2).
All terms of that sum are nonnegative, and this particular
term is $2+k_j\ge B(e_{j+1})\ge B(N)$.
This proves (IW.2) for every cutoff, not only a subsequence.
\end{proof}

\begin{corollary}[AC-K28-det: determinant identities do not remove the obstruction]
\label{AC-K28-det}
Let $f$ be as above and set
$\Lambda=\mathbb Z_p[[T_1,T_2]]$.
The perfect complex in degrees one and two
\[
 C=[\,\Lambda^2\xrightarrow{(x,y)\mapsto f(T_1)x}\Lambda\,]
                                                        \tag{IW.7}
\]
has $H^1(C)=\Lambda e_2$, $H^2(C)=\Lambda/(f)$, and the
inverse-determinant generator $z=f(T_1)e_2$ under the
generic trivialization.  In particular,
\[
 \operatorname{char}_\Lambda(H^1(C)/\Lambda z)
                   =\operatorname{char}_\Lambda H^2(C)=(f).
                                                        \tag{IW.8}
\]
Under the derived specialization
$T_1\mapsto(1+p)^e-1$, $T_2\mapsto0$, the first
cohomology is free of rank one, the second is finite,
and the specialized positive determinant index is
$v_p(f((1+p)^e-1))$.
Consequently perfectness, rank one, finiteness at all
weights, this characteristic-ideal identity, $\mu=0$,
and fixed one-variable Weierstrass degree two do not
imply any prescribed upper bound on the cumulative
positive indices, even on arbitrarily large sparse cutoffs.
\end{corollary}
\begin{proof}
The ring $\Lambda$ is a domain and $f\ne0$, so the
kernel of the displayed map is exactly $\Lambda e_2$,
and its cokernel is $\Lambda/(f)$.  Over its fraction
field, cancellation of the first coordinate and the
target in this two-term complex sends an inverse-
determinant basis to $f e_2$.  Equivalently the
two-term free resolution of $\Lambda/(f)$ contributes
the inverse-determinant lattice factor $f$; this is
the same sign convention as AC-K26.  Quotienting
$\Lambda e_2$ by $\Lambda f e_2$ proves (IW.8).
Derived specialization is calculated by specializing
this free complex itself.  By (IW.1) its resulting
kernel is $\mathbb Z_p e_2$ and its cokernel is
$\mathbb Z_p/(f((1+p)^e-1))$.
The claimed index and every finiteness statement follow.
Apply AC-K28 with $B$ larger than any proposed bound.
There is no assertion here that $C$ is a number-field
cohomology complex or that $z$ is an elliptic-unit class.
\end{proof}

\paragraph{Implication for the actual remaining gate.}
The arithmetic determinant generators of AC-K26 are highly
specific.  AC-K28 does not alter their identities or the
equivalence in AC-K27.  It proves why an argument for that
gate must use an additional property of those generators
which controls their proximity to the integer weight points.
An arithmetic height bound, an explicit integral cochain
estimate, or a regulator lower bound could provide such
information; none follows from the algebraic assertions
listed in AC-K28-det alone.  These examples also explain
why passing to sparse cutoffs is not by itself an algebraic
repair.  Their arbitrary $p$-adic coefficients, with no
height constraint, are an unrestricted input of the test,
not a statement about the coefficients furnished by arithmetic.

\paragraph{Dictionary scope.}
The complexes and inverse determinants lie in the same
perfect-complex and determinant-line dictionary already used
in AC-K26.  The finite quotients at specialization are
explicit modules over $\mathbb Z_p$, not motives or
constructed geometric Selmer groups.  No new claim about
a Grothendieck archive text is made.
